% No preamble. Prefix ahc-. Actual fixed-endpoint identity; uniform estimates conditional.
\section{The actual heat integral and a sufficient comparison estimate}
\label{ahc-section}

The Cartan identity also puts the missing comparison into
heat-kernel form without changing the lattice operators.
Use \(G(t),V_a(t)\) from \eqref{cdmom-probe-matrices}.
Let \(\operatorname{Tr}_{0,S}\) be the diagonal color trace
over noncenter scalar sites in block zero and
\(\operatorname{Tr}_{0,E}\) that over current bonds with
initial site in block zero. Define the formal coefficient
\begin{equation}
 h_{M,K}(s)=[t^2]\left\{
   \operatorname{Tr}_{0,E}e^{-sV_a(t)}
                  -2\operatorname{Tr}_{0,S}e^{-sG(t)}
                         \right\}.
 \label{ahc-h-definition}
\end{equation}
Every derivative here is evaluated at the flat positive
operators and on the polynomial stationary jets. No
heat evolution at a nonzero unbounded background is
used. Put
\[
 b_c(M,K)=[t^2]\sum_{c_-=0}
           \operatorname{tr}_{\mathfrak g}(\log P(t))_{cc}.
\]
Notebook 32 bounds \(|b_c|\) uniformly in \(M,K\).

\begin{proposition}[Exact actual heat representation]
\label{ahc-identity}
In each joint infinite system at fixed \(K\),
\begin{equation}
 \beta(D_{M,K})=b_c(M,K)
                   +\int_0^\infty h_{M,K}(s)\,\frac{ds}{s}.
 \label{ahc-exact}
\end{equation}
The integral is absolutely convergent for fixed \(K\).
The bounds establishing this statement may depend on \(K\).
\end{proposition}

\begin{proof}
First use bounded compactly supported output directions.
For a positive flat matrix family \(P(t)\), the local
functional calculus gives
\[
 [t^2](\log P(t))_{ii}
    =-\int_0^\infty [t^2](e^{-sP(t)})_{ii}\,\frac{ds}{s}.
\]
One can obtain this by differentiating the difference
of the proper-time formulas for \(P(t)\) and \(P(0)\);
their scalar reference term cancels.
To justify the same formula on polynomial jets, write
\(P(t)=P_0+tP_1+t^2P_2+O(t^3)\). Its heat coefficient is
\begin{align}
 [t^2]e^{-sP(t)}
 ={}&-\int_0^s e^{-(s-u)P_0}P_2e^{-uP_0}\,du \notag\\
 &+\int_{0<u<v<s}
 e^{-(s-v)P_0}P_1e^{-(v-u)P_0}P_1e^{-uP_0}\,du\,dv .
 \label{ahc-duhamel}
\end{align}
The \(P_j\) kernels are polynomially growing at worst,
with exponential localization in their relative
coordinates. This follows by inserting the affine
source jets in the established rooted derivative kernels.

The fixed-\(K\) flat gap and exponential matrix bound
give exponentially weighted semigroup estimates.
Indeed bounded truncated distance conjugation changes
\(P_0\) in operator norm by less than half its gap if
the exponent is small enough. The variation-of-constants
bound then gives
\(\|e^{-sP_0}\|_{\rm weighted}\le e^{-c_Ks}\).
Spending smaller spatial exponents absorbs the
polynomial factors in \eqref{ahc-duhamel}, after
anchoring at the diagonal block. Its trace is bounded by
\[
 C_K(s+s^2)e^{-c'_Ks},\qquad s>0,
\]
with constants uniform in volume and, where the
upstream bounds are, preceding depth. For small \(s\)
the same estimate follows directly from the one and
two insertion formulas. These bounds justify the
proper-time integral, polynomial-jet extension and
local infinite-volume passage by domination.
Substituting the two log coefficients in
\eqref{cdmom-exact-coefficient} proves
\eqref{ahc-exact}, with signs \(+1,-2\) in
\eqref{ahc-h-definition}.
\end{proof}

\subsection{A uniform short-time bound for the actual based ghost}
\label{ahc-ghost-short-section}

One part of the short-time comparison can be proved directly.
It uses the actual based scalar domain and has constants
independent of the preceding depth and endpoint.

\begin{proposition}[Actual ghost ultraviolet bound]
\label{ahc-ghost-short}
In four dimensions put
\[
 h^g_K(s)=-2[t^2]\operatorname{Tr}_{0,S}e^{-sG(t)}.
\]
For every permitted \(K\), on the affine probe already
identified with the actual stationary jets,
\begin{equation}
 0\le h^g_K(s)\le C_H\{(8s)^4+6(8s)^3+7(8s)^2+8s\},
                     \qquad s>0.
 \label{ahc-ghost-poisson}
\end{equation}
Consequently
\begin{equation}
 \int_0^1 |h^g_K(s)|\,\frac{ds}{s}\le2280C_H.
 \label{ahc-ghost-uv}
\end{equation}
Equivalently, the based-ghost part of the rescaled
short-time integral over \(0<u<K^{-2}\) is bounded
by this same constant.
\end{proposition}

\begin{proof}
On the complexified charge-\(q\) space, the actual based
operator is the covariant nearest-neighbor Laplacian
with Dirichlet values at every \(K\)-center. Hence
\[
 G_q(t)=8I-A_q(t),
\]
where \(A_q\) is the covariant adjacency operator on
noncenter sites. The diagonal is still \(8\), also
at sites adjacent to a removed center: it comes from
the full link norm of \(D_b\), with the center value
set to zero.

The diagonal of \(A_q(t)^n\) is a sum over length-\(n\)
closed nearest-neighbor walks avoiding the centers.
On the commuting affine field each walk has phase
\[
       \exp\{iqt\,{\cal A}_{\mu\nu}(\gamma)/K^2\},
\]
where \({\cal A}_{\mu\nu}(\gamma)\) is its oriented
algebraic area in the selected coordinate plane.
The contribution from a constant potential cancels
on a closed walk. Taking the origin of the area
coordinate at the starting site gives the elementary
bound \(|{\cal A}_{\mu\nu}(\gamma)|\le n^2\).
Thus the block trace coefficient is exactly
\begin{equation}
 h^g_K(s)=\frac{C_H}{K^4}e^{-8s}
   \sum_{\substack{z\ {\rm noncenter}\\z\ {\rm in\ block}\ 0}}
   \sum_{n\ge0}\frac{s^n}{n!}
   \sum_{\substack{\gamma:z\to z,\ |\gamma|=n\\
                    \gamma\ {\rm avoids\ centers}}}
                         {\cal A}_{\mu\nu}(\gamma)^2.
 \label{ahc-ghost-walk}
\end{equation}
All sums converge absolutely at each fixed \(s\);
the factorial absorbs the factor \(n^4\). This
also justifies extracting the formal quadratic
coefficient without imposing periodic magnetic flux.
Every term is nonnegative.

There are at most \(K^4\) starting sites and at most
\(8^n\) walks from each. It follows that
\[
 h^g_K(s)\le C_H e^{-8s}
                   \sum_{n\ge0}\frac{(8s)^n n^4}{n!}.
\]
The fourth moment of this exponential generating
series is \(x^4+6x^3+7x^2+x\) after multiplication
by \(e^{-x}\); it follows by applying
\((x\,d/dx)^4\) to \(e^x\). This proves
\eqref{ahc-ghost-poisson}. Integrating its four
monomials from zero to one with measure \(ds/s\)
gives
\[
 C_H\left(\frac{8^4}{4}
             +\frac{6\,8^3}{3}+\frac{7\,8^2}{2}+8\right)
                         =2280C_H.
\]
\end{proof}

This bound is uniform in \(M\) because the actual
stationary ghost jets on this probe have already
been identified with the same affine covariant
Laplacian for every \(S_M\). No estimate for a
separate massive inverse is used. The polynomial
upper bound in \eqref{ahc-ghost-poisson} is not a
useful large-time bound and does not settle the
long-time comparison.

\subsection{A quantitative target in physical heat time}
\label{ahc-target-section}

Set \(J_{M,K}(u)=h_{M,K}(K^2u)\), and
\(\alpha_H=11C_H/(48\pi^2)\).
The continuum calculation above has reference heat
coefficient
\[
 J_K^{\rm ref}(u)=\alpha_H e^{-u}
                         {\bf1}_{\{u\ge K^{-2}\}}.
\]
Thus the following is an exact difference identity,
not yet a bound:
\begin{equation}
 \beta(D_{M,K})-\beta_{\rm ref}(K^{-1},1)
 =b_c(M,K)+\int_0^\infty
          \{J_{M,K}(u)-J_K^{\rm ref}(u)\}\,\frac{du}{u}.
 \label{ahc-difference}
\end{equation}
It fixes the actual operator, regulator, sign, density,
and rescaling in a proposed comparison.

\begin{lemma}[Sufficient bounds on three heat-time intervals]
\label{ahc-band-lemma}
Suppose, for all required \(M,K\) and structural systems,
there are constants \(C,\delta>0\), independent of them,
such that
\begin{align}
 &\int_0^{K^{-2}}|J_{M,K}(u)|\,\frac{du}{u}\le C,
                                      \label{ahc-uv-bound}\\
 &|J_{M,K}(u)-\alpha_H|
       \le C\{(K^2u)^{-\delta}+u^\delta\},
          \quad K^{-2}\le u\le1,       \label{ahc-mid-bound}\\
 &\int_1^\infty|J_{M,K}(u)|\,\frac{du}{u}\le C.
                                      \label{ahc-ir-bound}
\end{align}
Then, with the already uniform coarse bound \(C_c\),
\begin{equation}
 \left|\beta(D_{M,K})
       -\frac{11C_H}{24\pi^2}\log K\right|
           \le C_c+2C+\frac{2C}{\delta}.
 \label{ahc-log-conclusion}
\end{equation}
\end{lemma}

\begin{proof}
Integrate the exact identity \eqref{ahc-exact} after
the time rescaling. The constant \(\alpha_H\) on the
middle interval integrates to \(2\alpha_H\log K\).
Each error in \eqref{ahc-mid-bound} integrates to
at most \(C/\delta\). The other intervals and the
coarse term give the remaining terms.
\end{proof}

The full bounds \eqref{ahc-uv-bound}--\eqref{ahc-ir-bound}
are not asserted here for the actual family uniformly in
\(K\). Proposition~\ref{ahc-ghost-short} supplies the
ghost part of the first bound; its vector part remains.
The fixed-endpoint heat representation by itself gives
only constants depending on \(K\). The middle interval asks for a
continuum coefficient on heat lengths between the
fine spacing and the final block size. The last
interval is where the soft pinned-ghost modes are
particularly relevant: the scalar gap alone may be
of order \(K^{-4}\) in fine units, so rescaling by
\(K^2\) does not give uniform damping. The combined
coefficient must be estimated, or a weaker signed
integral comparison must be proved. Bounding the
separate inverse operators by a physical mass is
not an available shortcut.

For \(M>1\), \(V_a\) contains the actual \(S_M\)
Hessian. Uniformity over that family is part of the
conditions, not a consequence of computing the bare
Wilson reference. A bounded comparison only at \(M=1\)
would not by itself give the varying-depth margin
needed by the finite drift-transfer theorem.

